% Authored for a11y-stem's development set (dev-v2). Original text, no third-party material.
% Licence: MIT. Author: a11y-stem contributors.
\documentclass[11pt]{article}
% Computer Modern text fonts, as pdfLaTeX writes them; XeTeX defaults to Latin Modern.
\usepackage[OT1]{fontenc}
\renewcommand{\rmdefault}{cmr}
\renewcommand{\ttdefault}{cmtt}
\usepackage{amsmath,amssymb}
\usepackage[margin=1in]{geometry}
\setlength{\parindent}{0pt}\setlength{\parskip}{8pt}
\begin{document}
\begin{center}{\Large\bfseries Determinants, Matrices and Systems: A Reference Sheet}\end{center}
\section*{Determinants}
A determinant is a single number computed from a square array. The examples below show how it is written and evaluated.
The determinant of a two by two matrix is the product of the diagonal entries minus the product of the others:
\label{eq000}\[
  \det A = \begin{vmatrix} 11 & -9 \\ 1 & -9 \end{vmatrix} = -90
\]
Expanding along the first row gives the determinant of a three by three matrix:
\label{eq001}\[
  \det B = \begin{vmatrix} 9 & 1 & 11 \\ 2 & 13 & -7 \\ 5 & 2 & -1 \end{vmatrix} = -695
\]
Scaling entries scales the terms of the determinant:
\label{eq002}\[
  \begin{vmatrix} 5a & b \\ 3c & d \end{vmatrix} = 5ad - 3bc
\]
The characteristic condition below is a quadratic equation in x:
\label{eq003}\[
  \begin{vmatrix} x - 3 & 11 \\ 1 & x + 3 \end{vmatrix} = 0
\]
Cramer's rule writes each unknown as a ratio of determinants:
\label{eq004}\[
  x = \frac{\begin{vmatrix} 1 & -3 \\ 5 & 3 \end{vmatrix}}{\begin{vmatrix} 13 & -3 \\ -15 & 3 \end{vmatrix}}, \quad y = \frac{\begin{vmatrix} 13 & 1 \\ -15 & 5 \end{vmatrix}}{\begin{vmatrix} 13 & -3 \\ -15 & 3 \end{vmatrix}}
\]
Vertical bars around an array also denote the determinant:
\label{eq005}\[
  \left| \begin{array}{cc} 11 & 7 \\ 15 & 9 \end{array} \right| = -6
\]
A row with a single non-zero entry makes the expansion short:
\label{eq006}\[
  \Delta = \begin{vmatrix} 3 & -9 & 3 \\ 0 & 5 & 0 \\ 3 & 5 & 7 \end{vmatrix} = 60
\]
The determinant of a triangular matrix is the product of its diagonal:
\label{eq007}\[
  \det C = \begin{vmatrix} -5 & -15 & 9 \\ 0 & 5 & 9 \\ 0 & 0 & 5 \end{vmatrix} = -125
\]
The determinant of a two by two matrix is the product of the diagonal entries minus the product of the others:
\label{eq008}\[
  \det A = \begin{vmatrix} 7 & -1 \\ 11 & 3 \end{vmatrix} = 32
\]
Expanding along the first row gives the determinant of a three by three matrix:
\label{eq009}\[
  \det B = \begin{vmatrix} 2 & 1 & 7 \\ 7 & 1 & 7 \\ -3 & -7 & -9 \end{vmatrix} = -200
\]
Scaling entries scales the terms of the determinant:
\label{eq010}\[
  \begin{vmatrix} 9a & b \\ 3c & d \end{vmatrix} = 9ad - 3bc
\]
The characteristic condition below is a quadratic equation in x:
\label{eq011}\[
  \begin{vmatrix} x - 7 & 9 \\ 1 & x + 7 \end{vmatrix} = 0
\]
Cramer's rule writes each unknown as a ratio of determinants:
\label{eq012}\[
  x = \frac{\begin{vmatrix} 2 & 11 \\ 1 & -3 \end{vmatrix}}{\begin{vmatrix} 15 & 11 \\ -9 & -3 \end{vmatrix}}, \quad y = \frac{\begin{vmatrix} 15 & 2 \\ -9 & 1 \end{vmatrix}}{\begin{vmatrix} 15 & 11 \\ -9 & -3 \end{vmatrix}}
\]
\section*{Matrices}
Matrices record data, transformations and systems of equations. Brackets and parentheses are both common.
Transposing a matrix swaps its rows and columns:
\label{eq013}\[
  A = \begin{bmatrix} -7 & -7 & -4 \\ -1 & 11 & 7 \end{bmatrix}, \qquad A^{T} = \begin{bmatrix} -7 & -1 \\ -7 & 11 \\ -4 & 7 \end{bmatrix}
\]
A matrix times a column vector gives another column vector:
\label{eq014}\[
  \begin{pmatrix} 2 & -7 \\ 1 & -9 \end{pmatrix}\begin{pmatrix} 4 \\ 0 \end{pmatrix} = \begin{pmatrix} 8 \\ 4 \end{pmatrix}
\]
The inverse of a two by two matrix uses the determinant and swaps entries:
\label{eq015}\[
  A^{-1} = \frac{1}{15} \begin{bmatrix} 3 & -1 \\ 9 & 2 \end{bmatrix}
\]
The augmented matrix of a system keeps the right-hand side beside the coefficients:
\label{eq016}\[
  \left[ \begin{array}{cc|c} -5 & 0 & 1 \\ -4 & -7 & 2 \end{array} \right]
\]
A three by three system gives an augmented matrix with one extra column:
\label{eq017}\[
  \left[ \begin{array}{ccc|c} -7 & 7 & 9 & 9 \\ 7 & 7 & 7 & 3 \\ 7 & 2 & -1 & 18 \end{array} \right]
\]
Row operations turn the left block into the identity and read off the solution:
\label{eq018}\[
  \left[ \begin{array}{cc|c} 3 & 3 & 9 \\ 9 & 1 & 19 \end{array} \right] \longrightarrow \left[ \begin{array}{cc|c} 1 & 0 & 2 \\ 0 & 1 & 1 \end{array} \right]
\]
Scalar multiples of matrices are added entry by entry:
\label{eq019}\[
  3\begin{bmatrix} -15 & -2 \\ 2 & 2 \end{bmatrix} + 7\begin{bmatrix} -11 & 11 \\ -5 & -4 \end{bmatrix} = \begin{bmatrix} -122 & 71 \\ -29 & -22 \end{bmatrix}
\]
A vector in three dimensions is written as a column:
\label{eq020}\[
  \mathbf{v} = \begin{pmatrix} 1 \\ -5 \\ 4 \end{pmatrix}
\]
A scalar multiple of the identity matrix has the same number down the diagonal:
\label{eq021}\[
  5I = \begin{bmatrix} 5 & 0 \\ 0 & 5 \end{bmatrix}
\]
Transposing a matrix swaps its rows and columns:
\label{eq022}\[
  A = \begin{bmatrix} 2 & 11 & 13 \\ -7 & 9 & -15 \end{bmatrix}, \qquad A^{T} = \begin{bmatrix} 2 & -7 \\ 11 & 9 \\ 13 & -15 \end{bmatrix}
\]
A matrix times a column vector gives another column vector:
\label{eq023}\[
  \begin{pmatrix} 7 & 3 \\ -15 & 11 \end{pmatrix}\begin{pmatrix} 4 \\ -1 \end{pmatrix} = \begin{pmatrix} 25 \\ -71 \end{pmatrix}
\]
The inverse of a two by two matrix uses the determinant and swaps entries:
\label{eq024}\[
  A^{-1} = \frac{1}{-59} \begin{bmatrix} 7 & -2 \\ 2 & -9 \end{bmatrix}
\]
The augmented matrix of a system keeps the right-hand side beside the coefficients:
\label{eq025}\[
  \left[ \begin{array}{cc|c} 3 & -15 & 15 \\ 3 & -15 & 1 \end{array} \right]
\]
A three by three system gives an augmented matrix with one extra column:
\label{eq026}\[
  \left[ \begin{array}{ccc|c} 9 & 1 & -9 & 4 \\ 15 & -5 & 9 & 7 \\ 3 & -3 & 13 & 14 \end{array} \right]
\]
Row operations turn the left block into the identity and read off the solution:
\label{eq027}\[
  \left[ \begin{array}{cc|c} 5 & 5 & 0 \\ -1 & 1 & -4 \end{array} \right] \longrightarrow \left[ \begin{array}{cc|c} 1 & 0 & 2 \\ 0 & 1 & -2 \end{array} \right]
\]
\section*{Systems of Equations}
A brace on the left groups equations that must hold together. Each system below has exactly one solution.
Solve the system of equations by elimination:
\label{eq028}\[
  \left\{ \begin{array}{l} -3x + 7y = 44 \\ 7x + 9y = 24 \end{array} \right.
\]
Find every pair of numbers that satisfies the system:
\label{eq029}\[
  \left\{ \begin{array}{l} -5x + 9y = -19 \\ 11x + 5y = 17 \end{array} \right.
\]
Use substitution to solve the following system:
\label{eq030}\[
  \left\{ \begin{array}{l} 3x + 11z = 8 \\ -5y - 9z = -24 \\ 11x - 7y - 2z = -34 \end{array} \right.
\]
Determine the values of the unknowns in the system:
\label{eq031}\[
  \left\{ \begin{array}{l} -3x + 3y = 6 \\ 7x + 7y = -28 \end{array} \right.
\]
Solve the system of equations by elimination:
\label{eq032}\[
  \left\{ \begin{array}{l} 7x + 11y = -29 \\ 7x - 2y = -3 \end{array} \right.
\]
Find every pair of numbers that satisfies the system:
\label{eq033}\[
  \left\{ \begin{array}{l} -9x + 7y - 3z = 10 \\ 9x + 2y - 9z = 11 \\ 9x + 3z = -3 \end{array} \right.
\]
Use substitution to solve the following system:
\label{eq034}\[
  \left\{ \begin{array}{l} 7x - 2y = 27 \\ 11x - 7y = 54 \end{array} \right.
\]
Determine the values of the unknowns in the system:
\label{eq035}\[
  \left\{ \begin{array}{l} -9x + 7y = 1 \\ 11x + 5y = 53 \end{array} \right.
\]
Solve the system of equations by elimination:
\label{eq036}\[
  \left\{ \begin{array}{l} 5x - 11y + 11z = -8 \\ 3x - 7y = -20 \\ 2x + 2z = 14 \end{array} \right.
\]
Find every pair of numbers that satisfies the system:
\label{eq037}\[
  \left\{ \begin{array}{l} -5x + 2y = 20 \\ 9x - 9y = -63 \end{array} \right.
\]
\section*{Worked Derivations}
Several lines of an argument can share one display, aligned on the equals sign.
Isolating the unknown takes three short steps:
\label{eq038}\begin{align*} 11x + 9 &= 31 \\ 11x &= 22 \\ x &= 2 \end{align*}
Expanding a square by hand shows where the middle term comes from:
\label{eq039}\begin{align*} (x + 3)^2 &= (x + 3)(x + 3) \\ &= x^2 + 3x + 3x + 9 \\ &= x^2 + 6x + 9 \end{align*}
Completing the square solves the quadratic in four steps:
\label{eq040}\begin{align*} x^2 + 14x + 5 &= 0 \\ (x + 7)^2 &= 44 \\ x + 7 &= \pm\sqrt{44} \\ x &= -7 \pm \sqrt{44} \end{align*}
A difference of squares factors at once:
\label{eq041}\begin{align*} x^2 - 121 &= x^2 - 11^2 \\ &= (x - 11)(x + 11) \end{align*}
The power rule is applied term by term:
\label{eq042}\begin{align*} \frac{d}{dx}\left(9x^3 + 3x\right) &= 9 \cdot 3x^2 + 3 \\ &= 27x^2 + 3 \end{align*}
Adding the two equations of a system removes y:
\label{eq043}\begin{align*} 2x + y &= 4 \\ x - y &= -1 \\ 3x &= 3 \\ x &= 1 \end{align*}
Matching bases turns an exponential equation into a linear one:
\label{eq044}\begin{align*} 2^{x + 6} &= 512 \\ 2^{x + 6} &= 2^{9} \\ x + 6 &= 9 \\ x &= 3 \end{align*}
A sum splits into a multiple of a known sum and a constant:
\label{eq045}\begin{align*} \sum_{k=1}^{10} (2k + 3) &= 2\sum_{k=1}^{10} k + 3 \cdot 10 \\ &= 2 \cdot 55 + 30 \\ &= 140 \end{align*}
Isolating the unknown takes three short steps:
\label{eq046}\begin{align*} 5x + 3 &= 13 \\ 5x &= 10 \\ x &= 2 \end{align*}
\end{document}
